%% notes-tex/025-graph-reg-sweep/sections/3-training.tex

\section{What the prior does to training\secstatus{todo}}
\label{sec:training}

\notesec{Support contraction happens at every $\lambda$, and saturates early}
Without a penalty, layer-1 attention has nothing to do with the graphs: edge recall at
degree is 0.006 at every diagnostic epoch, the value of attention unrelated to the
adjacency (Figure~\ref{fig:sweep}c). The smallest weight tested already moves it to
0.782 at epoch 29 (0.759 at epoch 9), $10^{-4}$ to 0.866, $10^{-3}$ to 0.906, and from
$10^{-2}$ up it sits at 0.914 to 0.915 with no further change to $\lambda = 1$. The
divergence tells the same story with more resolution (Figure~\ref{fig:sweep}d): from
1{,}540 at $10^{-5}$ through 571 and 337 to about 300 at $10^{-2}$, $10^{-1}$ and 1
(300, 297, 298), a floor the penalty cannot push below at any weight tested. Both
diagnostics are flat from $10^{-2}$ to 1 while the fixed-epoch accuracy still rises from
0.438 to 0.447. How close the attention sits to the graph is set by $\lambda = 10^{-2}$;
what the two larger weights add is not more contraction.

\notesec{The gradient budget}
The gradient probe (Figure~\ref{fig:sweep}e) measures the norm of each loss term's
gradient on one batch before clipping. The point-loss gradient at epoch 0 is a property
of the seed, not the arm: 7.4, 32.3 and 2.8 for seeds 1, 2 and 3 in every arm that
shares the initialization (the mask arm, with a different layer-1 kernel, reads 5.7,
32.9 and 2.7). The penalty's gradient scales with $\lambda$ as it must, 1.1 at
$10^{-3}$, 10.9 at $10^{-2}$, 109 at $10^{-1}$ and 1{,}097 at 1 in the seed mean at
epoch 0, and grows over training while the point-loss gradient settles to 1.0 to 4.3:
at epoch 20 the ratio is 1.1 to 1.5 at $10^{-3}$, 15 to 24 at $10^{-2}$, 86 to 146 at
$10^{-1}$ and 680 to 1{,}550 at $\lambda = 1$. The two best arms of the ladder train
under a total gradient that is almost entirely the penalty. The trainer clips the total
gradient norm at 10, so at $\lambda \geq 10^{-1}$ every update is rescaled by two to
three orders of magnitude before the optimizer sees it.

Hypothesis (untested): AdamW normalizes each parameter's update by its running gradient
scale, so a uniform rescaling of the whole gradient leaves the update nearly unchanged,
and the penalty's gradient is concentrated in the parameters that shape layer-1
attention (the layer-1 query and key projections and, through them, the gene table).
Under that reading the arms at $\lambda \geq 10^{-1}$ train the rest of the network on
the point loss much as the no-penalty arm does, which is consistent with their training
Pearson matching it (0.625 and 0.624 against 0.624), while layer 1 is held on the graph
for the whole run. The mechanism is not measured here; a per-parameter-group norm on
the probe batch would measure it.

\notesec{The prior is not a fit-reducing regularizer at the weights that help}
Training Pearson at epoch 29 (Figure~\ref{fig:sweep}h) is 0.624 with no penalty, 0.636
at $10^{-3}$, 0.613 at $10^{-2}$, 0.625 at $10^{-1}$ and 0.624 at 1, so the arms that
generalize best at epoch 29 fit the training triples as well as the arm that generalizes
worst. The two smallest weights fit the training set better than any other arm, 0.657
at $10^{-5}$ and 0.663 at $10^{-4}$, and generalize no better than no penalty. The
validation point loss (Figure~\ref{fig:sweep}i) reaches its minimum near epoch 13 in
every arm (0.80 to 0.82) and then rises; at epoch 29 it is 0.867 with no penalty and
0.869 under the mask against 0.826 and 0.825 at $\lambda = 10^{-1}$ and 1. The prior at
high weight does not stop the model from fitting the training triples; it slows the
growth of held-out error after the point-loss minimum.

Hypothesis (untested): the late decline of the unpenalized model is layer-1 attention
drifting away from any fixed structure as the gene table specializes to the training
triples, and a prior strong enough to hold layer 1 in place removes that degree of
freedom without touching the rest. The mask removes the same freedom in a different
way, by fixing the support and leaving the weights free, and it declines exactly as no
penalty does, so if this hypothesis is right the part of layer 1 that drifts is the
weighting within the support, not the support.
